% Informe técnico de socialización (español). Todas las cifras proceden del release con control de
% divulgación (release/public o, mientras no exista, release/synthetic con marca de agua).
\documentclass[11pt]{article}
\newcommand{\DocLang}{es}
\input{papers/shared/preamble.tex}
\input{papers/shared/authors.tex}
\input{papers/p2-lot-quality/shared/flow.tex}
\usepackage{tabularx}
\usepackage{fancyhdr}
\newcolumntype{Y}{>{\raggedright\arraybackslash}X}
\setlength{\headheight}{14pt}
\pagestyle{fancy}
\fancyhf{}
\fancyhead[L]{\small\textcolor{vgrey}{Informe técnico VA-MENGOC-BC}}
\fancyhead[R]{\small\textcolor{vgrey}{DICEI, Instituto Finlay de Vacunas}}
\fancyfoot[C]{\small\thepage}
\renewcommand{\headrulewidth}{0.4pt}
\hypersetup{pdftitle={Informe tecnico: farmacovigilancia y calidad de lotes de VA-MENGOC-BC},
  pdfauthor={Richard Alejandro Matos Arderi; Abel Ponce Gonzalez}}

\newcommand{\Mensaje}[1]{\item #1}

\begin{document}

\begin{titlepage}
\SyntheticNotice
\vspace*{2cm}
{\small\textcolor{vblue}{\textbf{INFORME TÉCNICO}}\par}
\medskip
{\LARGE\bfseries Evidencia reproducible sobre la vacuna VA-MENGOC-BC: seguridad poscomercialización y
calidad de los lotes liberados, \VStudyFirstYear--\VStudyLastYear\par}
\bigskip
{\large Informe para el grupo de expertos y la dirección del Instituto Finlay de Vacunas\par}
\vspace{1.5cm}
\AuthorBlock
\vfill
{\small\noindent
\begin{tabularx}{\linewidth}{@{}lY@{}}
Versión del código & \VCodeVersion{} (semilla \VSeed)\\
Origen de los datos & \ifVSynthetic Release sintético (datos simulados, no reales)\else Release
público con control de divulgación\fi\\
Tamaño mínimo de celda & \VMinCell{} (supresión primaria y complementaria)\\
Fecha de extracción & \VExtractionDate\\
Custodio de la base de datos & \VCustodian\\
Fecha de compilación & \today\\
\end{tabularx}\par}
\end{titlepage}

\tableofcontents
\clearpage

\section*{Mensajes principales}
\addcontentsline{toc}{section}{Mensajes principales}
\begin{itemize}[leftmargin=*,itemsep=3pt]
\Mensaje{Se analizaron \PONReports{} notificaciones de eventos supuestamente atribuibles a la
  vacunación o inmunización (ESAVI) sin duplicados; \PONTarget{} (\POPctTarget) incluyeron
  VA-MENGOC-BC.}
\Mensaje{La tasa global de notificación fue de \POPooledRate{} por \num{100000} dosis (\CI{}
  \POPooledRateCI), \PORateVsExpected{} la tasa de referencia del programa, y \POTrendWord{} entre
  \PORateFirstYear{} y \PORateLastYear{} (razón de tasas anual \POTrendRR, \CI{} \POTrendRRCI).}
\Mensaje{Con el criterio preespecificado, \POSignalsPrimaryAllOtherVaccines{} tipos de evento se
  notificaron de forma desproporcionada (\POSignalList); \PONSignalsRobustPhrase{} se mantuvo en los
  \PONDesigns{} diseños de comparación. Las señales no descritas como esperadas
  (\POSignalEmergingList) requieren revisión clínica caso a caso.}
\Mensaje{Se enlazaron \PTNLinked{} de \PTNTargetReports{} notificaciones (\PTLinkRate) con un lote
  liberado. \PTConformitySentence.}
\Mensaje{De \PTNPrimaryTests{} asociaciones preespecificadas entre atributos de los lotes y
  reacciones, \PTNSignificantFDR{} superaron la corrección por comparaciones múltiples
  (\PTSignificantList). \PTNegControlSentence.}
\Mensaje{Ninguna persona es identificable en los resultados: el release aplica supresión de celdas
  menores que \VMinCell{} y protege los totales frente a la recuperación por diferencia.}
\end{itemize}

\section{Objetivo y alcance}
El proyecto genera, con un flujo de trabajo reproducible, la evidencia de dos artículos científicos
sobre VA-MENGOC-BC, la vacuna cubana de vesículas de membrana externa contra los meningococos de los
serogrupos B y~C:
\begin{enumerate}[leftmargin=*]
\item \textbf{Artículo 1} (\emph{Drug Safety}): perfil de seguridad notificado en el sistema nacional
  de vigilancia pasiva, con tasas de notificación, análisis de desproporcionalidad según READUS-PV,
  fenotipos latentes de eventos notificados conjuntamente y factores asociados a la hospitalización.
\item \textbf{Artículo 2} (\emph{Vaccine}): calidad de los lotes liberados y asociación entre la
  variación dentro de especificación de sus atributos y el perfil de reactogenicidad notificado,
  según STROBE y RECORD-PE.
\end{enumerate}
Este informe resume los métodos, los resultados y las decisiones pendientes para el grupo de
expertos. No sustituye a los manuscritos, que contienen la descripción completa y la bibliografía.

\section{Datos y gobernanza}
\subsection{Fuentes}
Se utilizaron \VNFiles{} ficheros anuales de notificaciones de ESAVI (\VStudyFirstYear--%
\VStudyLastYear) del sistema nacional de vigilancia y el registro de liberación de \PTNLots{} lotes
producidos entre \PTLotFirstYear{} y \PTLotLastYear{} (\PTNNationalLots{} nacionales y
\PTNExportLots{} de exportación). De \VReportsRaw{} registros leídos se eliminaron
\VExactDuplicates{} duplicados exactos y se fusionaron \VKeyDuplicatesMerged{} duplicados por clave;
\VReportsInWindow{} notificaciones quedaron dentro de la ventana de estudio.

\subsection{Protección de los datos}
Los datos personales de salud y la información industrial nunca salen de la estación controlada. El
tratamiento sigue los principios ALCOA+, la seudonimización con clave secreta
\citep{ENISA2019,MHRA2018}, el control estadístico de la divulgación \citep{Hundepool2012} y la
Ley~149/2022 de Protección de Datos Personales \citep{Ley149}:
\begin{itemize}[leftmargin=*]
\item nombres y direcciones se usan solo para derivar un seudónimo HMAC-SHA256 y se descartan; la
  fecha de nacimiento se transforma en edad y el texto libre se categoriza y se elimina;
\item cada ejecución registra la versión del código, la configuración, la semilla y la huella
  SHA-256 de cada fichero de entrada;
\item el release público suprime los recuentos menores que \VMinCell{} (marca ``<\VMinCell'') y
  aplica supresión complementaria (marca [c]) cuando un recuento suprimido podría recuperarse por
  diferencia; los atributos de los lotes se publican solo como posición normalizada dentro de la
  ventana de especificación.
\end{itemize}

\subsection{Calidad de los datos}
La \tabref{tab:completeness} resume la completitud de los campos clave por año y la
\tabref{tab:plausibility}, las comprobaciones de plausibilidad. Los registros marcados se informan y
no se eliminan.

\begin{table}[htbp]
\centering\small
\caption{Completitud de los campos clave por año del fichero (\% de registros con el campo
presente).}
\label{tab:completeness}
\POTableCompleteness
\end{table}

\begin{table}[htbp]
\centering\small
\caption{Comprobaciones de plausibilidad: registros marcados en todos los años.}
\label{tab:plausibility}
\POTablePlausibility
\end{table}

\section{Flujo de procesamiento}
El pipeline es un paquete de Python con pruebas automáticas que se ejecuta con una sola orden
(\texttt{uv run vamengoc run}). La \tabref{tab:pipeline} describe sus ocho pasos.

\begin{table}[htbp]
\centering\small
\caption{Pasos del flujo de procesamiento y garantías que ofrece cada uno.}
\label{tab:pipeline}
\begin{tabularx}{\linewidth}{@{}>{\raggedright\arraybackslash}p{3.1cm}Y>{\raggedright\arraybackslash}p{4cm}@{}}
\toprule
Paso & Qué hace & Garantía\\
\midrule
1. Ingesta & Lee los ficheros tal como llegan; detecta cabeceras desplazadas y deriva del esquema
  entre años & Originales inmutables y trazables\\
2. Seudonimización & Seudónimo HMAC; elimina nombres, direcciones y texto libre & Minimización desde
  la primera línea\\
3. Normalización & Tipa fechas, edades, dosis, lotes y especificaciones con reglas explícitas &
  Sin limpieza manual\\
4. Calidad de datos & Completitud, conformidad y plausibilidad por campo y año & Calidad medida
  antes de analizar\\
5. Enlace con lotes & Une cada notificación con su lote (exacto, núcleo o aproximado) & Información
  industrial protegida\\
6. Análisis & Análisis preespecificados de ambos artículos con semillas fijas & Resultados
  reproducibles\\
7. Control de divulgación & Supresión primaria y complementaria; verificación independiente &
  Nadie identificable\\
8. Publicación & Tablas, figuras y macros para manuscritos, informes, web y presentaciones & Texto y
  resultados coherentes\\
\bottomrule
\end{tabularx}
\end{table}

\section{Resultados del artículo 1: seguridad poscomercialización}
\subsection{Notificaciones y tasas}
La mediana del tiempo entre la vacunación y la notificación fue de \PODelayMedian{} días. La
\figref{fig:rates} muestra la serie mensual de notificaciones y la tasa anual de VA-MENGOC-BC por
\num{100000} dosis, con la tasa de referencia del programa (\POExpectedRate{} por \num{100000}
dosis). Se detectaron \PONBreaksAllPhrase{} en la serie mensual de todas las vacunas
(\POBreaksAll).

\begin{figure}[htbp]
\centering
\includegraphics[width=\textwidth]{p1_rates}
\caption{(a) Notificaciones mensuales de ESAVI para todas las vacunas y para VA-MENGOC-BC, con los
puntos de cambio detectados por PELT. (b) Tasa anual de notificación de VA-MENGOC-BC por
\num{100000} dosis con intervalo de confianza exacto del 95\,\%.}
\label{fig:rates}
\end{figure}

\subsection{Desproporcionalidad}
Se evaluaron \PONEventsScreened{} tipos de evento con cuatro métodos (ROR, PRR, componente de
información y EBGM). El criterio preespecificado fue IC\textsubscript{025}${}>0$ con al menos
\PODispMinReports{} notificaciones. La \tabref{tab:designs} muestra el número de señales en cada
diseño de comparación y la \figref{fig:forest}, las razones de odds de notificación.

\begin{table}[htbp]
\centering\small
\caption{Número de señales de desproporcionalidad en los diseños de comparación preespecificados.}
\label{tab:designs}
\POTableDesigns
\end{table}

\begin{figure}[htbp]
\centering
\includegraphics{p1_forest}
\caption{Razones de odds de notificación (IC 95\,\%, escala logarítmica) de los eventos con
VA-MENGOC-BC frente al resto de las vacunas y en el diseño de comparador activo en lactantes.}
\label{fig:forest}
\end{figure}

\subsection{Fenotipos latentes y hospitalización}
El análisis de clases latentes seleccionó \POLCABestK{} clases, estables entre remuestreos (mediana
del índice de Rand ajustado \POLCAAri; \figref{fig:lca}). En la validación temporal
(entrenamiento \POTrainYears, prueba \POTestYears), el mejor modelo de hospitalización alcanzó un
área bajo la curva ROC de \POBestAuroc, una discriminación \PODiscriminationWord.
\POHospInterpretation.

\begin{figure}[htbp]
\centering
\includegraphics[width=\textwidth]{p1_lca}
\caption{Perfiles de las clases latentes: probabilidad de cada evento en cada clase y tamaño
relativo de las clases.}
\label{fig:lca}
\end{figure}

\section{Resultados del artículo 2: calidad de los lotes y reactogenicidad}
\subsection{Enlace notificación--lote}
La \figref{fig:flow} resume el enlace y las exclusiones. Se analizaron \PTNAnalysed{} notificaciones
en \PTNAnalysedLots{} lotes (mediana de \PTReportsPerLot{} notificaciones por lote).

\begin{figure}[htbp]
\centering
\LinkageFlow
\caption{Enlace de las notificaciones de VA-MENGOC-BC con los lotes liberados.}
\label{fig:flow}
\end{figure}

\subsection{Calidad del proceso}
\PTConformitySentence. Se detectaron \PTNChangepoints{} puntos de cambio en la secuencia de
producción (\PTChangepointList). La \figref{fig:capability} muestra la capacidad de proceso por
atributo y período; los atributos con Ppk inferior a~1 fueron: \PTPpkBelowOneList. Los lotes
nacionales y de exportación fueron equivalentes en \PTNEquivalent{} de \PTNAttributes{} atributos
(margen del \PTEqMarginPct\,\% de la ventana).

\begin{figure}[htbp]
\centering
\includegraphics[width=\textwidth]{p2_capability}
\caption{Índice de capacidad de proceso Ppk por atributo y período de producción, con intervalo de
confianza bootstrap del 95\,\%.}
\label{fig:capability}
\end{figure}

\subsection{Asociación entre atributos de los lotes y reacciones}
El diseño de solo casos compara el perfil de reacciones entre notificaciones de lotes con distintos
valores de los atributos, con ecuaciones de estimación generalizadas que tienen en cuenta la
correlación dentro del lote \citep{Liang1986} y la corrección de Benjamini--Hochberg
\citep{Benjamini1995}. La \tabref{tab:gee} presenta las estimaciones y la \figref{fig:p2forest}, su
representación gráfica. \PTSensSentence. \PTNegControlSentence.

\begin{table}[htbp]
\centering\small
\caption{Razones de odds por desviación estándar (IC 95\,\%) de las asociaciones preespecificadas:
estimaciones crudas y ajustadas por edad, sexo, dosis, año y coadministración; $q$, valor de $p$
ajustado por Benjamini--Hochberg.}
\label{tab:gee}
\PTTableGEE
\end{table}

\begin{figure}[htbp]
\centering
\includegraphics[width=\textwidth]{p2_forest}
\caption{Razones de odds por desviación estándar de cada atributo para las reacciones principales,
con la banda del efecto mínimo detectable.}
\label{fig:p2forest}
\end{figure}

\section{Limitaciones}
\begin{itemize}[leftmargin=*]
\item La vigilancia pasiva no permite medir la infranotificación ni la notificación estimulada; la
  desproporcionalidad genera hipótesis y no estima riesgos.
\item Los eventos no se validaron frente a las definiciones de caso de Brighton; las señales no
  descritas requieren revisión caso a caso.
\item No se dispone de dosis por edad, número de dosis ni provincia, ni de dosis distribuidas por
  lote: no pueden calcularse tasas estratificadas ni incidencias por lote.
\item Los ensayos de liberación describen el lote en el momento de liberarse, no tras el
  almacenamiento y la distribución.
\end{itemize}

\section{Decisiones que se solicitan al grupo de expertos}
\begin{enumerate}[leftmargin=*]
\item Confirmar el significado de las siglas de las vacunas y de los códigos de provincia.
\item Fijar la fecha que define el año analítico (vacunación, notificación o año del fichero).
\item Validar la clasificación de gravedad de los eventos y decidir si se codifican con MedDRA o
  CIE-10.
\item Autorizar, si procede, la publicación de algún atributo de calidad en unidades absolutas.
\item Aportar las dosis administradas en \VStudyLastYear{} y confirmar la tasa de referencia del
  programa.
\item Organizar la revisión caso a caso de las señales no descritas por el comité nacional.
\item Completar la aprobación ética, la financiación y las contribuciones de los autores antes del
  envío.
\end{enumerate}

\section{Reproducibilidad y material disponible}
\begin{itemize}[leftmargin=*]
\item Código, configuración y release público: \url{https://github.com/Pol4720/va-mengoc-bc-ml}.
\item Aplicación web interactiva y presentaciones:
  \url{https://pol4720.github.io/va-mengoc-bc-ml/}.
\item Regeneración completa en la estación controlada: \texttt{uv run vamengoc run} seguido de
  \texttt{uv run vamengoc release verify public}; los manuscritos, este informe y la web se
  actualizan a partir del nuevo release sin editar ninguna cifra a mano.
\end{itemize}

\bibliographystyle{vancouver}
\bibliography{papers/shared/bib/pubmed,papers/shared/bib/manual}

\end{document}
